%%%%%
%
\section{Fonction exponentielle polynomiale}

\itemstyle{Op\'erateur~:} \textsf{Poly\_exp}
\index{Poly\_exp}

%%%%%OBJET%%%%
\itemstyle{Objet~:}
Application d'une fonction exponentielle polynomiale (de degr\'e 2).

%%%%%ARGUMENTS%%%%
\bigskip
%%%%%ENTREES%%%%
\begin{tabarg}
\multicolumn{3}{l}{\titretab {Entr\'ees}}\\
\hline
$X$ & Vecteur, {\tt float} & Op\'erande (unit\'e quelconque)\\
$C_0$ & Coeff., {\tt double} & Coefficient du terme de degr\'e 0.\\
$C_1$ & Coeff., {\tt double} & Coefficient du terme de degr\'e 1.\\
$C_2$ & Coeff., {\tt double} & Coefficient pour l'exposant du terme de degr\'e 2.\\
$C_3$ & Coeff., {\tt double} & Coefficient pour l'exposant du terme de degr\'e 1.\\
$C_4$ & Coeff., {\tt double} & Coefficient pour l'exposant du terme de degr\'e 0.\\
$Unit\_{res}$ & Coeff., {\tt String} & Unit\'e du r\'esultat.\\
\end{tabarg}

\medskip
%%%%%SORTIES%%%%
\begin{tabarg}
\multicolumn{3}{l}{\titretab {Sorties}}\\
\hline
$Z$ & Vecteur, {\tt float} & R\'esultat (unit\'e indiqu\'ee par $Unit\_{res}$)\\
\end{tabarg}


%%%%%FORMULE%%%%
\itemstyle{Formule}

%% Poly_exp
\[ Z = C_0 + C_1 . e^{\displaystyle{(C_2 + C_3 . X + C_4 . X^2)}} \]


%%%%%COMMENTAIRE%%%%
\itemstyle{Commentaire~:}
On calcule d'abord une variable interm\'ediaire Y qui est une fonction polynomiale de X (degr\'e 2).
Puis, on calcule le r\'esultat Z , fonction exponentielle de Y.

%%%%%AUTEUR%%%%
\itemstyle{Auteur~:}
CNRM/GMEI/TRAMM

%%%%%REFERENCE%%%%
\itemstyle{R\'ef\'erences~:}
